The counts of one-centimeter and three-centimeter sticks have the same parity because their sum is 100. If both counts are even, pair equal sticks, choose one pair as two opposite sides, and split the remaining pairs equally between the other two sides. If both counts are odd, take one stick of each length and break the three-centimeter stick into lengths one and two. Make one side of length two from the original unit stick and the unit fragment, and the opposite side from the length-two fragment. The remaining counts are both even, so divide matching pairs between the other two sides. All sides have positive length and at most one stick was broken. 
\begin{asy}
texpreamble("\usepackage{amsmath,amssymb}");
size(220);
defaultpen(fontsize(9pt));
draw((0,0)--(5,0)--(5,2)--(0,2)--cycle);label("equal paired totals",(2.5,0),(0,-1));label("$2$",(5,1),(1,0));draw((0,1)--(.15,1));label("$1+1$",(0,1),(-1,0));
\end{asy}
